\chapter{INTRODUCTION}
\label{chap:introduction}

%─────────────────────────────────────────────────────────────────────────────
\section{Background and Motivation}
%─────────────────────────────────────────────────────────────────────────────

Classical robot systems are typically task-specific pipelines in which perception,
planning, and control are engineered independently for a narrow operating scenario.
Such systems can perform well in controlled laboratory conditions, but they are often
inflexible in real-world environments where objects, layouts, and goals change
frequently. Even small distribution shifts --- such as unseen object geometry,
different lighting, or a reformulated command --- can cause significant degradation in
performance and require expensive redesign.

A second limitation is poor generalization. Many robot models learn behavior tied to
the training distribution and fail to transfer reliably to unseen tasks. This limits
practical deployment in settings where robots must adapt quickly to new goals without
full retraining. At the same time, human-robot interaction remains a bottleneck:
existing systems often rely on rigid APIs, manually crafted state machines, or
specialized operators rather than natural language instructions.

Recent Vision-Language Models (VLMs) provide strong multimodal reasoning and
instruction-following capability~\cite{openai2023gpt4v,liu2023llava},
but using a VLM alone is not sufficient for robust robot control. Practical embodied
systems require explicit integration between visual perception, language understanding,
and action generation. In many current implementations, this integration remains weak
or ad hoc, resulting in fragmented pipelines and unstable behavior under real-time
constraints.

This thesis therefore targets a general-purpose Vision-Language-Action (VLA) system
for robot manipulation. The proposed approach integrates three cooperating modules:
(1) a VLM core for language-conditioned scene understanding, (2) a Vision Expert for
task-relevant visual feature extraction, and (3) an Action Expert for structured
action generation. The goal is to enable robots to follow natural-language
instructions, improve generalization to unseen tasks, and maintain real-time
performance under practical compute and data constraints.


%─────────────────────────────────────────────────────────────────────────────
\section{Problem Identification}
%─────────────────────────────────────────────────────────────────────────────

Based on the identified limitations of prior work, the current problem is framed by
the following gaps.

\paragraph{Gap 1: Task-Specific and Inflexible Robot Pipelines.}
Classical robot systems are strongly engineered for predefined tasks and often break
under changing objects, layouts, or goals. Their dependence on handcrafted modules
reduces scalability across diverse real-world conditions.

\paragraph{Gap 2: Weak Generalization to Unseen Tasks.}
Many existing policies perform well only on seen task templates and fail to adapt to
new task compositions or novel object arrangements, limiting practical utility.

\paragraph{Gap 3: Limited Natural-Language Instruction Following.}
Most deployed systems still rely on rigid command interfaces instead of natural
language, preventing intuitive interaction by non-expert users and increasing
integration overhead.

\paragraph{Gap 4: Incomplete Vision-Language-Action Integration.}
Although vision and language models have advanced rapidly, many robot systems still
lack efficient end-to-end integration between perception, language grounding, and
action generation. This motivates a unified VLA architecture.

Building on these identified gaps, this thesis poses the following primary research
question:

\begin{quote}
\textit{How can a Vision-Language-Action architecture integrate a VLM, a Vision
Expert, and an Action Expert to enable natural-language robot manipulation while
improving generalization to unseen tasks under practical real-time constraints?}
\end{quote}

The following secondary questions are also addressed:
\begin{enumerate}[label=\arabic*)]
    \item How should a VLA system be modularized so that VLM reasoning, visual
    feature extraction, and action generation operate coherently in one pipeline?
    \item Which design choices are most effective for balancing action expressiveness
    with real-time latency and memory constraints?
    \item To what extent does the proposed system generalize to unseen robot
    manipulation tasks compared with seen-task performance?
\end{enumerate}


%─────────────────────────────────────────────────────────────────────────────
\section{Objectives}
%─────────────────────────────────────────────────────────────────────────────

The overarching objective of this research is to design, implement, and evaluate a
general-purpose Vision-Language-Action system for robot manipulation. The specific
objectives are as follows:

\begin{enumerate}[label=\arabic*)]
    \item To develop a Vision-Language-Action (VLA) system for robot manipulation.
    \item To integrate a VLM core with a Vision Expert and an Action Expert in a
    unified architecture.
    \item To enable robots to follow natural-language instructions robustly.
    \item To improve generalization to unseen tasks and environments.
    \item To evaluate system performance on representative robot manipulation tasks.
\end{enumerate}


%─────────────────────────────────────────────────────────────────────────────
\section{Scope and Limitations}
%─────────────────────────────────────────────────────────────────────────────

The scope of this research is delimited as follows. The target domain is robot arm
manipulation in indoor workspaces using RGB visual input and natural-language user
commands. The system focuses on instruction following and action generation for
manipulation tasks, not full autonomous navigation. The architecture is restricted to
three integrated modules (VLM, Vision Expert, Action Expert), and evaluation is
performed on representative manipulation tasks under practical compute constraints.
Depth sensing, large-scale multi-robot coordination, and formal safety verification
are outside the current scope.


%─────────────────────────────────────────────────────────────────────────────
\section{Contributions}
%─────────────────────────────────────────────────────────────────────────────

The principal contributions of this thesis are as follows:

\begin{enumerate}[label=\arabic*)]
    \item \textbf{A practical Vision-Language-Action architecture} that unifies
    multimodal understanding and action generation for robot manipulation.

    \item \textbf{An integrated three-module design} that combines a VLM, a Vision
    Expert, and an Action Expert with explicit interfaces for stable inference.

    \item \textbf{A natural-language robot control pipeline} enabling users to issue
    plain-English manipulation instructions without low-level programming.

    \item \textbf{A generalization-oriented training and evaluation setup} targeting
    performance on unseen manipulation tasks rather than only seen templates.

    \item \textbf{A task-based experimental benchmark} that evaluates instruction
    following, action validity, generalization, and latency on representative robot
    manipulation categories.
\end{enumerate}


%─────────────────────────────────────────────────────────────────────────────
\section{Thesis Organisation}
%─────────────────────────────────────────────────────────────────────────────

The remainder of this thesis is structured as follows.
Chapter~\ref{chap:litreview} reviews related theories and prior work on multimodal
robot learning, natural-language instruction following, and action modeling for
manipulation.
Chapter~\ref{chap:methodology} presents the proposed VLA system design, including
module integration (VLM + Vision Expert + Action Expert), design constraints, and
implementation strategy.
Chapter~\ref{chap:experiments} reports experimental results on manipulation tasks,
including instruction-following performance, generalization to unseen tasks, and
real-time latency analysis.
Chapter~\ref{chap:conclusion} summarizes key findings, limitations, and future
directions for general-purpose intelligent robot systems.